\section{Background and Related Work}

\subsection{Traditional vs.\ Agentic Scientific Workflows}

High Performance Computing workflows consist of multi-stage applications that
produce and manipulate data in multiple discrete phases to produce an end
result or insight. A phase typically corresponds to a computational task
(e.g., a simulation, analysis, or visualization step) that consumes data
produced by previous phases and generates data for subsequent ones. For
instance, a workflow meant to discover new alloys could involve preparing
alloy structures in one phase, running Density Functional Theory calculations
on each, and then aggregating and visualizing the data. Sometimes these phases
may be cyclical, with phases repeating in a set order, where the outputs of
one series of phases become the inputs to the next. Common high-impact HPC
workflows include protein and alloy discovery, astronomical simulation, and
fluid dynamics workflows. The use of these workflows to simulate, analyze, and
derive insights from data has led to significant advances in multiple
scientific fields, including the discovery of new proteins for medical
applications, analysis and identification of cosmological phenomena such as
supernovae in massive telescope data streams, and rapid speed-ups of
prototyping for fluid-dependent systems in fields such as aerospace
engineering.

\subsection{Prior Solutions for Memory Bottlenecks}

Despite their advantages for scientific progress, HPC workflows face a key
challenge: their operations on massive datasets lead to an imbalance between
the efficiency of compute hardware and memory bandwidth. Workflows often
involve multi-GB datasets that can balloon quickly as simulation steps produce
more data from their initial inputs (e.g., 100,000 atoms having their energy
interactions for every 1~ns step simulated by a protein folding program). This
data can easily make an HPC workflow's data ``out-of-core'' for part or all of
a workflow, meaning the data is too large to fit in RAM all at once. This
results in data movement costs, and the infamous ``memory wall'': hardware
FLOPs have scaled $60000\times$ in the past 20 years, while DRAM bandwidth has
only scaled $100\times$ and interconnect bandwidth $20\times$ in the same
timespan. While some HPC sites have addressed this by scaling their DRAM, the
financial and energy cost of this scaling makes it unsustainable.

In recent work, the solution to this issue has come in the form of software
distributed shared memory (DSM) systems that present tiered memory and storage
resources as an ``infinite memory pool'', often referred to as a Deep Memory
and Storage Hierarchy (DMSH). This system relies on the predictable
phase-based nature of HPC workflows in order to handle both the organization
of data across the memory and storage hierarchy and ensure coherency without
infeasible overhead. Information about phase intent and accesses allows data to
be staged proactively, placed across the hierarchy intelligently, and
prefetched while other computation is running to ensure no idle computational
resources. For example, if a simulation phase is known to consume a particular
trajectory, a DMSH data staging system can begin loading that data from
storage while the preceding analysis phase is still executing, hiding most of
the I/O latency behind computation. Intent and expected accesses are
transmitted to the system via API.

\subsection{Workload Description}

\paragraph{AtomAgents.}
AtomAgents is a multi-agent materials discovery framework built on AutoGen
that combines LLM reasoning with physics-based simulation tools such as
LAMMPS. Specialized agents collaborate to plan, execute, and refine scientific
workflows by invoking simulation tools through function-calling interfaces.
Each tool invocation represents a workflow phase whose required simulation
inputs, intermediate datasets, and even successor phases are determined only
after the agents complete their reasoning.
